\section{Relation to prior research}
The nearest marketing work is RAMP, which already connects client-specific memory, planning and verification for audience curation~\citep{p250811120}. Its combined interventions do not isolate a memory-only effect, and our constrained account-action task is not its audience-selection benchmark. We therefore make no cross-paper numerical leaderboard claim. The retained critical review compares its methods, controls and limited human assessment with this study.

Agent-memory systems and evaluations already address persistent facts, experience, temporal updates and retrieval. Our prior review covers Hindsight, Zep, Mem0, LongMemEval and domain-specific memory evaluation~\citep{p251212818,p250113956,p250419413,p241010813,p260726072}. Wiki compilation and compression work, including WiCER and CMC, further motivates strong generic representations and fixed-reader controls~\citep{p260507068,p260932658}. Case-based reasoning and truth maintenance are older foundations for connecting evidence to actions and maintaining justified beliefs~\citep{aamodt1994,doyle1979}. Our contribution is narrower: an auditable context-budget comparison connected to feasible next-action decisions, with a control that can reproduce the specialized selector using query words alone.

Model-style agent benchmarks provide useful evaluation principles. $\tau$-bench checks domain-policy tasks against database outcomes and distinguishes average success from repeated-trial reliability \citep{taubench2024}. AppWorld allows multiple valid solutions, checks collateral changes and uses executable validation solutions; its held-out-application challenge is stronger generalization evidence than a new generator seed \citep{appworld2024}. $\tau^2$-bench composes verified subtasks and separates interaction modes to diagnose failure mechanisms \citep{tau2bench2025}. We adopt explicit constraints, acceptable alternatives and independent executable references. Our single-step, clean-record task does not reproduce their interactive environments or justify calling our scores general model capability.
